\section{Finite graph regularity at scri: a normal-bootstrap lemma}
\label{sec:constraint-normal-bootstrap}

This section records the precise normal-bootstrap argument separately from
weak existence. It is a linear regularity statement, not a sharp
$H^6\longrightarrow H^6$ geometric evolution theorem. The hypotheses below
include the exact normal reduction; a formal indicial polynomial alone would
not suffice.

\paragraph{Normal variables and exact structure.}
Work in a collar where $r$ is bounded away from zero. Put
\[
 v=(\tau,\zeta),\qquad N=\Omega\partial_\Omega v,
 \qquad\gamma=-\partial_r\Omega=\frac r{aS}>0.
\]
The derivative-reduction identities recover $N$ without a new differentiability
assumption:
\[
 N_\tau=Q_n/G_n-2\tau,\qquad
 N_{\zeta_j}=Z_{nj}/G_n-\zeta_j .
\]
All angular differentiations below are rotation Lie derivatives on the
appropriate scalar/vector/tensor fields; they span tangential differentiations
uniformly on the collar. They commute with radial dilation and with the
rotationally covariant equations.

Eliminate $H,M$ using the exact equations
\begin{align}
 H[T,z]&=\frac{2\Omega}{A}D_0T-2\Omega^2\partial_i z_i
             +2d_\theta T,\\
 M_i[T,z]&=\frac\Omega A(D_0z_i+2lz_i)
             -\partial_iT-cz_i,
 \qquad D_0=\partial_t-B\cdot\partial .
\end{align}
Substituting these expressions in the exact Hamiltonian and momentum
subsidiary equations gives four second-order equations with common principal
part $D_0^2-A^2\Delta$. Its radial second-order coefficient is
$B_r^2-A^2=-\Omega^2$ exactly. After $T=\Omega^2\tau$,
$z=\Omega\zeta$, multiplication by $\Omega^2$, and introduction of $N$,
the equations have the structural form
\begin{equation}
 -\gamma^2\Omega^3\partial_\Omega N+
 L(\Omega,n)N+K(\Omega,n)v+
 2B_r\gamma\Omega\partial_tN
 +\Omega\mathcal T_1v+\Omega^2\mathcal T_2v=0 .
 \label{eq:normal-eight-structural}
\end{equation}
The matrix coefficients are smooth and bounded up to scri. The operators
$\mathcal T_j$ involve only time and angular differentiations and have order
at most $j$; zero-order terms can be included in them. There are no hidden
negative powers or angular derivatives of $N$. This follows from the exact
second-order equations, not from freezing their principal symbols. The
unscaled vector equations contain $T/\Omega^3$, but the displayed balanced
weights turn that term into a bounded coefficient of $v$ in
\eqref{eq:normal-eight-structural}.

At scri the normal matrix is
\[
 L_0=2A_{\scri}^2lJ_{gg},
\]
where the scalar/normal block of $J_{gg}$ is
\[
 \begin{pmatrix}
 -l-d_\theta/2&-\kappa_1/2\\
 l+\kappa_1\kappa_2/2&2l/3-\kappa_1/2
 \end{pmatrix}
\]
and each transverse entry is $-\kappa_1/2$. The common wave metric makes
$-L_0$ strictly coercive. Thus, in a smaller collar, $L$ is invertible and
$-L$ remains coercive. Define
\[
 F=-L^{-1}K,
 \qquad R_g=N-Fv .
\]
The exact slow equation is
\begin{equation}
 (\Omega\partial_\Omega-F_0)v=R_g+\Omega F_1v,
 \qquad F=F_0+\Omega F_1 .
 \label{eq:normal-slow}
\end{equation}
Equation \eqref{eq:normal-eight-structural} gives a fast equation with leading
operator $\Omega^3\gamma^2\partial_\Omega-L$ on $R_g$. Its remaining terms
are sums of $\Omega$ times time/angular operators of order at most one and
$\Omega^2$ times operators of order at most two on $v,R_g$.
For a fixed-time estimate, the $\Omega\partial_tR_g$ term is placed on the
right. Differentiating $Fv$ produces $\Omega^2FN$ and smooth higher-order
terms, because $\partial_\Omega v=N/\Omega$; these introduce no unrecorded
normal derivative and no worse power.

\paragraph{Indicial data.}
For $a\kappa_1=4$, $\kappa_2=-99/100$, the scalar/normal block of $F_0$ is
\[
 \begin{pmatrix}58/31&-150/31\\24/155&242/31\end{pmatrix};
\]
its eigenvalues are $2,238/31$. The transverse blocks are $2I$.
The root $2$ is semisimple with multiplicity three. Spectral projectors are
smooth rotationally covariant functions of $n$. There is no integer
resonance between the two roots. The leading root-two coefficients satisfy
\[
 (h/l,m_n/l,\tau,\zeta_n)
   =(-2/25,302/75,1,-2/75)\tau,
 \qquad m_A=-l\zeta_A.
\]

\paragraph{Two one-dimensional estimates used below.}
First, the fast operator has a fixed-time inverse on every fixed power
weight. For
\[
 \Omega^3a(\Omega)\partial_\Omega u+D(\Omega)u=f,
 \qquad a>0,\quad\operatorname{sym}_HD\ge d_*H>0,
\]
integration by parts against $\Omega^p u^*H$ absorbs the coefficient-derivative
error, bounded by $C_p\Omega_0^2\|u\|_p^2$, in a sufficiently small collar.
The outer endpoint term is nonnegative. The inner growing homogeneous branch
is excluded by membership in a power-weighted $L^2$ space. Hence the inverse
is bounded in $L^2(\Omega^p d\Omega)$ for every fixed real $p$. The collar
may be chosen once for the finite set of weights used below. For constant
scalar $D$ the formula is
\[
 u(\Omega)=e^{D/(2\Omega^2)}
 \int_0^\Omega e^{-D/(2s^2)}f(s)\,s^{-3}\,ds,
\]
whose kernel mass is $1/D$. This fixes the sign of the admissible inverse.

Second, diagonalizing $F_0$ reduces
$(\Omega\partial_\Omega-F_0)v=f$ to scalar Euler equations.
Their Green operators are bounded between equal $b$-weights whenever the
weight differs from an indicial root, with constants controlled by the
inverse distance to that root. When improving a weight across an indicial
root, the corresponding $\Omega^\lambda$ term is extracted. This is obtained
directly by integrating
$\partial_\Omega(\Omega^{-\lambda}v)=\Omega^{-\lambda-1}f$ and applying the
one-dimensional weighted Hardy inequality, or equivalently Young's inequality
after $s=-\log\Omega$. The outer-collar trace is controlled by the local
$L^2$ norms of $v,N$ on a fixed overlap annulus. No independent data at scri
are introduced.

\paragraph{Finite normal-bootstrap statement.}
Let a reduction-compatible linear subsidiary solution have finite, continuous
$L^2$ \emph{joint} graph norms
\[
 \mathcal G_{40}(t)=
 \sum_{p+|I|\le40}
 \|\mathcal Z^I\partial_t^pU(t)\|_{L^2_H}.
\]
Here $\mathcal Z$ contains $\mathcal D=\Omega\partial_\Omega$ on the collar,
the three covariant rotations, and cutoff Cartesian derivatives spanning
the interior, with the overlapping cutoffs defined in the existence section.
All words in these generators are included.  The initial phase therefore
includes conormal regularity \emph{from the outset}; it is propagated by the
first-order commutator proof in that section.  Time derivatives are actual
generator derivatives, not freely prescribed corner jets.  This strengthens
the initial domain.  The earlier assertion that time/rotation-only graph
regularity supplies the conormal levels is withdrawn.
Assuming the exact structural reduction
\eqref{eq:normal-eight-structural}, it has an expansion
\begin{equation}
 v=\sum_{m=2}^7\Omega^m a_m(t,n)+R,
 \qquad \Omega^{-15/2}R\in H_b^6
 \label{eq:normal-finite-expansion}
\end{equation}
locally in the collar, with estimates by $C\mathcal G_{40}(t)$ and the fixed
overlap-annulus norms. The coefficient $a_2$ lies in the root-two kernel;
the integer coefficients $a_3,\ldots,a_7$ are fixed by recursion and the
time/angular derivatives of lower coefficients. The second algebraic branch
$\Omega^{238/31}$ is allowed in $R$.

\paragraph{Proof and derivative bookkeeping.}
Use $b$-measure $d\Omega/\Omega\,d\omega$, and let $\alpha$ denote the
normal weight, so that a weighted norm is the norm of
$\Omega^{-\alpha}u$. Ordinary collar $L^2$ is the starting $b$-weight
$\alpha=-1/2$. Both $v$ and $N$, hence also $R_g$, are bounded linear
functions of $U$ in this norm. The same holds after every finite joint
time/rotation/conormal word in the available budget.  The cutoff Cartesian
generators provide the corresponding ordinary estimates in the interior
and on the overlap annulus.

If the finite joint family of $v,R_g$ is controlled at weight $\alpha$ with
two additional time/rotation generators,
the fast right hand side is controlled at weight $\alpha+1$: its terms have
at least one factor of $\Omega$, and at most two time/angular differentiations.
The fixed-time fast estimate therefore controls $R_g$ at weight
$\alpha+1$.  This is applied \emph{simultaneously} to the required conormal
levels: the source $\Omega\partial_t\mathcal D^jR_g$ is known from the
previous weight at the same conormal level and one higher time level.
It is not used to construct an unknown conormal derivative.
The commutator of $\mathcal D^j$ with $\Omega^3\partial_\Omega$ contributes
$\Omega^2$ times the current $\mathcal D^jR_g$ and lower levels.  The current
term is absorbed into a small smooth change of the coercive fast matrix.
Thus the same fixed-time inverse applies at all these levels.
Equation \eqref{eq:normal-slow} then controls $v$ at the same
improved weight, modulo crossed indicial terms. It also controls $N=Fv+R_g$.
The other fields are reconstructed by
\begin{align}
 Q&=G(2\tau+N_\tau)+\Omega\nabla_{\rm tan}\tau,\\
 Z_{ij}&=G_i(\zeta_j+N_{\zeta_j})
                 +\Omega\partial_{{\rm tan},i}\zeta_j,\\
 h&=\frac{2\Omega}{A}\partial_t\tau+2wN_\tau
                  +(4w+2d_\theta)\tau-2\operatorname{tr}Z,\\
 m&=\frac\Omega A\partial_t\zeta+wN_\zeta
                  +(w-c+2l\Omega/A)\zeta-Q.
\end{align}
Thus the weight improvement returns the full field, not only $v$.

When a known polynomial is present, subtract its recursively determined
fast and slow polynomials before applying this estimate. Smooth coefficients
and time/angular differentiation preserve these finite expansions. At each
new integer order $m\ne2$, $(mI-F_0)$ is invertible; at $m=2$ its kernel is
the outgoing coefficient already identified. The coefficient is bounded by
the Euler Green estimate from the preceding weight and forcing norm.
There are no root-zero or root-one terms because neither is an indicial
root and the initial weight improvements exclude them. Eight improvements
\[
 -\tfrac12\longrightarrow\tfrac12\longrightarrow\cdots
 \longrightarrow\tfrac{15}{2}
\]
extract precisely the powers $\Omega^2,\ldots,\Omega^7$ and cost at most
sixteen time/angular generators. Endpoints never hit an indicial root;
$238/31-15/2=11/62>0$.

Six conormal derivatives are reserved in the starting joint norm and are
carried through the simultaneous weight induction.  The identity
$\mathcal Dv=Fv+R_g$ is compatible with this hierarchy, but is not used to
infer missing conormal norms.  Six angular differentiations and two spare
time differentiations give the conservative budget
$16+6+6+2=30<40$.  At each weight gain the remaining time/rotation budget
decreases by at most two while the reserved conormal levels remain present.
This removes the circular step that would result from putting
$\Omega\partial_t\mathcal DR_g$ on the right before knowing
$\mathcal DR_g$ at the preceding weight.  No infinite regularity or
unrecorded negative-weight commutator is used.

The use of fixed-time fast inversion is essential to this bookkeeping.
If instead the $\Omega\partial_tR_g$ term were kept on the left in a
space-time estimate, an initial half-weight trace would appear and would
require a separate estimate. It is not silently included here.

\paragraph{Ordinary finite regularity and scope.}
After factoring $v=\Omega^2v_{\rm nat}$, the remainder in
\eqref{eq:normal-finite-expansion} has $b$-weight $11/2$. For
$V=\Omega^{-11/2}R_{\rm nat}$,
\[
 \partial_\Omega^6R_{\rm nat}
 =\Omega^{-1/2}\prod_{j=0}^5
       (\Omega\partial_\Omega+11/2-j)V.
\]
The squared ordinary measure $d\Omega$ is exactly the $b$-measure
$d\Omega/\Omega$ on the right. Mixed derivatives are controlled with the
remaining angular budget. Consequently the natural constraint fields
$(H/\Omega^4,T/\Omega^4,M/\Omega^3,z/\Omega^3)$ have ordinary collar
$H^6$ regularity and finite traces. Their leading trace relations are the
root-two kernel relations above. The generic fractional natural exponent
$176/31$ need not be $H^7$; no infinite smoothness is asserted.

This is a nonsharp graph-domain embedding. If the subsidiary semigroup and
its maximal graph domain have been established separately, its strengthened
joint graph norms
propagate without derivative loss in that chosen topology, and the embedding
gives preservation of these natural falloffs. It does not imply an ordinary
$H^6$-data estimate for arbitrary data, does not replace the exact normal
reduction by its indicial polynomial, and does not establish the full
geometric/elliptic-gauge evolution.

